\FloatBarrier
\section{Coherence, Relaxation, and Pulse Errors}
\label{app:noise-leakage}

\label{sec:noise-leakage-robust-control}

During a drive-free wait, fluctuations of the spin splitting reduce
phase coherence, while energy exchange with the environment changes
spin populations. During a gate, voltage fluctuations can change the
intended rotation angle, and the pulse can populate states outside
the computational space. This appendix develops the corresponding
calculations behind Section~\ref{sec:platform-coherence}: Ramsey and
echo probabilities, population relaxation, and pulse-dependent angle
and leakage errors. The conditions for reducing each error depend on
the waiting configuration, noise, and pulse shape.

\subsection{Correcting and characterizing evolution during a wait}
\label{sec:idle-single-spin-characterization}

Calibrated $\pi/2$ and $\pi$ pulses prepare and analyze a spin
before and after a drive-free wait. The pulse conventions are those
of Section~\ref{sec:calibration-spin-rotations}. The detuning during the wait contains two physically different
parts,
\[
 \delta(t)=\delta_{\mathrm r}+\xi(t).
\]
The reproducible part $\delta_{\mathrm r}$ produces a Ramsey-fringe
phase drift and can be corrected by changing the reference frequency.
The random part $\xi(t)$ reduces the fringe contrast and defines the
phase-coherence time. An echo pulse cancels a constant or sufficiently
slow contribution during the wait and exposes the remaining phase
noise. A separate relaxation experiment measures the loss of
upper-state population. The resulting $T_2^*$, $T_{2,\mathrm{echo}}$
and $T_1$ belong to the chosen waiting configuration.

\subsubsection{Measuring phase coherence with Ramsey and echo}
\label{sec:measurement-free-spin-dynamics}
\label{sec:measurement-ramsey-echo}

The resonance scan in Section~\ref{sec:measurement-rabi} determines
$\omega_{\mathrm{res}}$. The calibrated rotations of
Section~\ref{sec:calibration-spin-rotations} use
$\omega=\omega_{\mathrm{res}}$, so their nominal detuning is
$\delta=\omega_{\mathrm{res}}-\omega=0$. Ramsey now tests what
happens to the spin phase during a drive-free wait. Finite precision
in the frequency scan or later drift may leave a residual mean
detuning $\delta_{\mathrm r}$. The spin-transition frequency
$\omega_q(t)=E_Z(t)/\hbar$ may also fluctuate, causing different
repetitions to accumulate different phases relative to the fixed
control reference. The first $\pi/2$ pulse creates a superposition,
and the second pulse converts this relative phase into a logical-basis
population. The phase of the resulting fringe reveals
$\delta_{\mathrm r}$, while the decay of its contrast defines the
phase-coherence time $T_2^*$.

Use the logical basis $\ket0=\ket{\downarrow}$,
$\ket1=\ket{\uparrow}$ and $Z=\ket0\bra0-\ket1\bra1$ from
Section~\ref{sec:single-qubit-gates}. Let
$Y=-i\ket0\bra1+i\ket1\bra0$ and
$R_y(\theta)=e^{-i\theta Y/2}$. All states and rotations below
are expressed in the frame rotating at the drive angular frequency
$\omega=\omega_{\mathrm{res}}$. The gates and the waiting evolution
must use this same frame. Turning off the transverse drive means
setting its envelope to $\Omega(t)=0$; it does not mean setting
$\omega=0$. The oscillator that generates the two pulses retains
the phase reference $\omega t$, so the spin can accumulate a
measurable phase relative to that reference while no transverse
field is applied.

Equation~\eqref{eq:calibration-single-spin-hamiltonian} specifies the pulses:
choose drive phases $\phi_1=-\pi/2$ and $\phi_2=+\pi/2$,
with the same amplitude $\Omega_0$ and duration $t_p$.
For $|\delta_{\mathrm r}|\ll\Omega_0$, the pulse Hamiltonians and
duration are
\[
 H_1\simeq\frac{\hbar\Omega_0}{2}Y,\qquad
 H_2\simeq-\frac{\hbar\Omega_0}{2}Y,\qquad
 t_p=\frac{\pi}{2\Omega_0}.
\]
They give $R_y(\pi/2)$ and $R_y(-\pi/2)$, respectively.
The phase difference of $\pi$ reverses the drive term;
the drive envelope is $\Omega(t)=0$ during the intervening wait,
while the oscillator phase reference continues at angular frequency
$\omega$.
Treating these pulses as short compared with the phase-evolution
and decoherence times gives the ideal Ramsey sequence
\begin{equation}
 \ket0\xrightarrow{R_y(\pi/2)}
 \frac{\ket0+\ket1}{\sqrt2}
 \xrightarrow{\text{wait }\tau}
 \frac{\ket0+\ket1}{\sqrt2}
 \xrightarrow{R_y(-\pi/2)}\ket0.
 \label{eq:measurement-ramsey-sequence}
\end{equation}
If a residual mean detuning $\delta_{\mathrm r}$ is present but
random fluctuations are absent, the drive-free wait adds the phase
$\delta_{\mathrm r}\tau$, and the inverse-pulse sequence predicts
\begin{equation}
 P_0(\tau)=\frac{1+\cos(\delta_{\mathrm r}\tau)}2.
 \label{eq:measurement-ramsey-ideal}
\end{equation}
This residual is the difference between the mean transition
frequency during the Ramsey acquisition and the previously chosen
reference frequency. It is therefore a correction to the nominally
resonant setting, not a second independent determination of
$\omega_{\mathrm{res}}$.

Use the Rabi measurement in Section~\ref{sec:measurement-rabi}
to determine $t_p$ at the chosen amplitude. In each trial,
prepare $\ket0$, apply the first pulse, wait $\tau$ with the
drive off, and apply the second $\pi/2$ pulse. Let $\phi_a$ denote
an additional phase applied to this second pulse relative to the
inverse pulse in Eq.~\eqref{eq:measurement-ramsey-sequence}.
Then use the single-spin readout of
Section~\ref{sec:calibration-single-spin-readout}. For each fixed
$\tau$, scan $\phi_a$ through a full period and repeat the complete
sequence at every phase to estimate $P_0(\tau,\phi_a)$. Keep the
field, waiting gate settings, drive frequency, amplitude and pulse
duration fixed. The probability-versus-phase curve is the Ramsey
fringe; its amplitude measures the reproducible coherence after
the wait. Repeat this phase scan for several values of $\tau$ to
obtain the coherence decay.

The energy difference between the two spin states may fluctuate
during the wait. Write
\[
 E_Z(t)=\overline E_Z+\varepsilon_Z(t),\qquad
 \frac{\overline E_Z}{\hbar}
       =\omega_{\mathrm{res}}+\delta_{\mathrm r},
\]
where $\overline E_Z$ is the mean splitting during the acquisition,
$\delta_{\mathrm r}$ is its residual detuning from the fixed control
reference $\omega_{\mathrm{res}}$, and $\varepsilon_Z(t)$ is a
zero-mean energy fluctuation. Define
\[
 \xi(t)=\frac{\varepsilon_Z(t)}{\hbar}
\]
only to express this energy noise in angular-frequency units. Model
$\xi(t)$ as a classical random function that can differ between
repetitions. Assume that the fluctuation changes the gap while
leaving $\ket0$ and $\ket1$
as the energy eigenstates, and neglect relaxation and pulse errors.
With the drive reference fixed at $\omega_{\mathrm{res}}$, the
waiting Hamiltonian in the same rotating frame is
\[
 H_{\mathrm{wait}}(t)
   =-\frac{\hbar}{2}
      \bigl[\delta_{\mathrm r}+\xi(t)\bigr]Z.
\]
All these Hamiltonians are proportional to $Z$, so they commute
at different times. The waiting evolution is therefore
\[
 U_{\xi}(\tau)
   =\exp\!\left\{\frac{iZ}{2}
                 \bigl[\delta_{\mathrm r}\tau+\varphi(\tau)\bigr]
          \right\},
 \qquad
 \varphi(\tau)=\int_0^\tau\xi(t)\,dt.
\]
For a given function $\xi(t)$ in one repetition, the
wait adds the phase $\delta_{\mathrm r}\tau+\varphi(\tau)$.
The analysis pulse with
phase offset $\phi_a$ is
\[
 R_{\mathrm a}(\phi_a)
 =\exp\!\left[-\frac{i\pi}{4}
       (-\sin\phi_a X-\cos\phi_a Y)\right].
\]
It reduces to $R_y(-\pi/2)$ when $\phi_a=0$. The conditional return
probability is
\[
 \begin{aligned}
 P_0(\tau,\phi_a\mid\xi)
   &=\left|\bra0 R_{\mathrm a}(\phi_a)U_{\xi}(\tau)
                         R_y(\pi/2)\ket0\right|^2\\
   &=\frac{1+\cos[\delta_{\mathrm r}\tau
                         +\varphi(\tau)-\phi_a]}{2}.
 \end{aligned}
\]
Repeated trials sample different fluctuations as well as the
quantum readout outcomes. Their counts estimate the average of
this probability. Below, $\langle\cdot\rangle$ denotes a
classical average over the fluctuations in those repetitions.
Using $\cos x=\operatorname{Re}e^{ix}$ gives
\begin{equation}
 \begin{aligned}
 P_0(\tau,\phi_a)
   &=\left\langle
       P_0(\tau,\phi_a\mid\xi)\right\rangle\\
   &=\frac12\left[1+
       \operatorname{Re}\{e^{i(\delta_{\mathrm r}\tau-\phi_a)}
                            W_R(\tau)\}\right]\\
   &=\frac12\left[1+|W_R(\tau)|
       \cos\bigl(\phi_a-\delta_{\mathrm r}\tau
                         -\arg W_R(\tau)\bigr)\right],\\
 W_R(\tau)&=\left\langle e^{i\varphi(\tau)}\right\rangle.
 \end{aligned}
 \label{eq:measurement-ramsey-noise}
\end{equation}
Thus a sinusoidal fit to the measured phase scan directly gives
the contrast $|W_R(\tau)|$ and the fringe phase
$\delta_{\mathrm r}\tau+\arg W_R(\tau)$. The contrast describes
how averaging the random phase over repeated trials reduces the
reproducible coherence. If the fluctuations have zero mean and no
systematic phase, the slope of the fringe phase versus $\tau$ gives
$\delta_{\mathrm r}$.

For example, suppose $\xi(t)$ is constant during each
repetition but Gaussian distributed between repetitions, with
zero mean and standard deviation $\sigma_\xi$. Then
$\varphi(\tau)=\xi\tau$, and averaging
$e^{i\xi\tau}$ over that Gaussian distribution gives
\begin{equation}
 W_R(\tau)=e^{-\sigma_\xi^2\tau^2/2}
           =e^{-(\tau/T_2^*)^2},\qquad
 T_2^*=\frac{\sqrt2}{\sigma_\xi}.
 \label{eq:measurement-quasistatic-dephasing}
\end{equation}
Thus Ramsey contrast can decay even when each individual spin
undergoes coherent evolution: different repetitions accumulate
different phases.

The measured contrasts at different waiting times give the
coherence envelope. A Gaussian fit yields $T_2^*$ under the
quasistatic Gaussian-noise model above. The fitted fringe phases
give the residual mean detuning and reveal any additional mean
phase shift not included in the zero-mean noise model.

To remove an approximately constant frequency error within a
repetition, insert a $\pi$ rotation about a transverse axis halfway
through the wait. It interchanges the two spin amplitudes, so they
accumulate the opposite relative phase during the second half.
For ideal pulses, the remaining noise phase is
\begin{equation}
 \varphi_E(\tau)=\int_0^{\tau/2}\xi(t)\,dt
                   -\int_{\tau/2}^{\tau}\xi(t)\,dt.
 \label{eq:measurement-echo-phase}
\end{equation}
The residual mean term contributes
$\delta_{\mathrm r}\tau/2-\delta_{\mathrm r}\tau/2=0$, and any
time-independent part of $\xi$ cancels in the same way.
This is a \emph{Hahn echo}; its return
contrast depends on $|\langle e^{i\varphi_E}\rangle|$ and is often
described by an echo decay time $T_{2,\mathrm{echo}}$. Fluctuations
that change during the sequence, spin relaxation, and imperfect
pulses still reduce the contrast.

For the echo experiment, repeat the same measurement with an
additional calibrated $\pi$ pulse halfway through the wait.
Vary the total waiting time and obtain the return-fringe contrast
at each value. Comparing Ramsey and echo contrasts tests how much
of the apparent dephasing is removed by refocusing slow frequency
fluctuations. These two sequences have been implemented with
electrically driven silicon spins~\cite{Kawakami2014}.

\subsubsection{Population relaxation during a wait}
\label{sec:measurement-relaxation}

After a resonant $\pi$ pulse, the electron is in the upper spin
state. With the drive off it can release energy to its environment
and return to the lower state, through the mechanism described in
Section~\ref{sec:calibration-single-spin-initialization}.
This loss of population does not follow from the unitary waiting
Hamiltonian $H_{\mathrm{wait}}$; it is described by transition
rates. We measure how long the upper-state population survives at
the chosen holding voltages. If upward and downward transitions have constant rates
$\gamma_\uparrow$ and $\gamma_\downarrow$, the population obeys
\begin{equation}
 \dot P_\uparrow=-\gamma_\downarrow P_\uparrow
                   +\gamma_\uparrow(1-P_\uparrow),\qquad
 P_\uparrow(\tau)=P_\uparrow^{\mathrm{eq}}
      +[P_\uparrow(0)-P_\uparrow^{\mathrm{eq}}]e^{-\tau/T_1},
 \label{eq:measurement-t1-decay}
\end{equation}
where $T_1=(\gamma_\downarrow+\gamma_\uparrow)^{-1}$ and
$P_\uparrow^{\mathrm{eq}}=\gamma_\uparrow/
(\gamma_\downarrow+\gamma_\uparrow)$. At low temperature the
upward rate is small and the upper-state population decays toward
zero.

To measure this decay, initialize down, apply the $\pi$ pulse of
Section~\ref{sec:calibration-spin-rotations}, turn off the drive,
wait at a fixed gate setting for time $\tau$, and apply the same
readout pulse after every wait. Repeat for several $\tau$ values
and fit the excited-state probability to
Eq.~\eqref{eq:measurement-t1-decay}. The fitted $T_1$ belongs to
the waiting configuration. Keeping the readout configuration and
its acquisition window fixed prevents their changes from being
confused with a change in the waiting-time decay.

\subsection{Reducing sensitivity to voltage fluctuations}

Consider exchange evolution that acts as a rotation about one
fixed logical axis with angle $\theta=\int J(t)\,dt/\hbar$.
Examples are a 1--2 pulse of the exchange-only qubit and an
exchange-dominated singlet--triplet $z$ rotation with its
gradient negligible during the pulse. Voltage fluctuations change
the confinement and the tunnel barrier, so nominally identical
pulses can have different exchange energies and accumulated phases.
For example, at aligned dot energies, $J=4t_c^2/(U-V_{12})$:
a fluctuating overlap already changes $J$ even if the mean dot
energies are held fixed.

\subsubsection{Connecting voltage noise to an error in the gate angle}
Let the actual voltages be $\mathbf V(t)+\delta\mathbf V$, where
$\mathbf V(t)$ is the programmed pulse and the offsets
$\delta V_g$ are nearly constant during one pulse but vary between
repetitions. For a pulse about one fixed axis, the random extra
energy produces an extra rotation angle,
$\delta\theta=\hbar^{-1}\int[J(\mathbf V+\delta\mathbf V)
-J(\mathbf V)]dt$. To first order,

\begin{equation}
 \delta\theta
   =\sum_g b_g\delta V_g,\qquad
 b_g=\frac1\hbar\int_0^\tau
        \left.\frac{\partial J}{\partial V_g}\right|_{\mathbf V(t)}dt.
 \label{eq:calibration-angle-noise}
\end{equation}
Here the programmed waveform $\mathbf V(t)$ is fixed, while the
offsets are random. With zero-mean offsets and covariance
$C_{gh}=\mathbb E[\delta V_g\delta V_h]$, their angle variance is
\[
 \mathbb E[(\delta\theta)^2]=\sum_{g,h}b_g C_{gh}b_h.
\]
For an input perpendicular to the rotation axis, the overlap with
the desired output is $\cos^2(\delta\theta/2)$.
Its mean loss is approximately
$\mathbb E[(\delta\theta)^2]/4$ for small fluctuations.
Thus reducing the voltage derivatives that enter $b_g$ directly
reduces this contribution to gate error.

For an explicit prediction of the averaged signal, take Gaussian
offsets. Then $\delta\theta$ is Gaussian with variance
$\sigma_\theta^2=\mathbf b^TC\mathbf b$. In a basis where the
rotation axis is $Z$, the error multiplies $\rho_{01}$ by
$e^{-i\delta\theta}$ in each trial. Averaging over repetitions gives
\begin{equation}
 \mathbb E[\rho_{01}']=e^{-\sigma_\theta^2/2}\rho_{01},
 \qquad
 \mathbb E[F_\perp]=\frac{1+e^{-\sigma_\theta^2/2}}2.
 \label{eq:calibration-angle-noise-channel}
\end{equation}
Here the intended rotation has been removed and $F_\perp$ is
the overlap for an input perpendicular to the axis.
The mean error over all pure states of this logical qubit is smaller,
$(1-e^{-\sigma_\theta^2/2})/3$, because states on the rotation
axis are unaffected. This gives a gate-error prediction from
the measured voltage response and a specified noise covariance,
rather than only a qualitative preference for a small derivative.

\subsubsection{Choosing a less sensitive exchange setting}
The derivative formula tells us what to reduce, but does not yet
choose a voltage setting. The symmetric two-dot model gives a
concrete choice. Holding
$t_c$ and $D=U-V_{12}$ fixed,
\begin{equation}
 J(\varepsilon)=\frac{4t_c^2D}{D^2-\varepsilon^2}
 =\frac{4t_c^2}{D}
   \left[1+\frac{\varepsilon^2}{D^2}
           +O\!\left(\frac{\varepsilon^4}{D^4}\right)\right].
 \label{eq:calibration-symmetric-exchange}
\end{equation}
At $\varepsilon=0$, a small detuning offset has no first-order
effect on $J$. One can therefore hold the dot energies aligned
while changing the barrier to control exchange. This is the
principle of symmetric exchange operation~\cite{Reed2016}.
Barrier fluctuations still change $t_c$, and hence $J$; the
suppressed derivative is specifically the detuning derivative.

Locate this setting experimentally by measuring $J$ on both sides
of the aligned-dot point. Use the measured response along the
actual compensated voltage path, since a voltage can also change
$t_c$ and $D$. Compare candidate pulses at the same target angle:
a smaller instantaneous derivative may be offset by a longer
gate duration.

\subsubsection{Estimating what remains after the first-order term vanishes}
At the symmetric point, detuning noise still changes exchange
through the quadratic term. Its residual effect can be estimated
with the same fixed-angle calculation. For a constant pulse with
$\theta_0=4t_c^2\tau/(\hbar D)$ and a small quasistatic detuning
offset $\delta\varepsilon$,
\[
 \delta\theta\simeq\theta_0\frac{(\delta\varepsilon)^2}{D^2}.
\]
If $\delta\varepsilon$ is zero-mean Gaussian with variance
$\sigma_\varepsilon^2$, then
$\mathbb E[(\delta\varepsilon)^4]=3\sigma_\varepsilon^4$.
Consequently,
\begin{equation}
 \mathbb E[\delta\theta]=
       \theta_0\frac{\sigma_\varepsilon^2}{D^2},\qquad
 \operatorname{Var}(\delta\theta)=
       2\theta_0^2\frac{\sigma_\varepsilon^4}{D^4}.
 \label{eq:calibration-quadratic-detuning-noise}
\end{equation}
The mean shift can be included in the angle calibration.
After that correction, the leading mean gate infidelity from
this noise alone is $\operatorname{Var}(\delta\theta)/6$.
It is fourth order in $\sigma_\varepsilon/D$, instead of the
second-order dependence produced by a nonzero linear derivative.
This improvement does not apply to a fluctuating barrier:
$\delta J/J\simeq2\delta t_c/t_c$ still gives a first-order
angle fluctuation. The working-point choice has removed first-order sensitivity to
detuning but not to barrier noise. Compare such settings using
pulses with the same requested angle.

\subsection{Avoiding population of excluded states}

A voltage pulse reshapes the orbital potential as well as changing
the intended spin interaction. An electron prepared in the lowest
orbital can then acquire amplitude in an excited orbital.
For a concrete example, let $\ket1$ contain an up spin in the
orbital ground state, and let $\ket\ell$ have the same spin in
the first excited orbital along $x$. We calculate this leakage
channel; the other logical state can have its own channels.

\subsubsection{Identifying the transition driven by the voltage pulse}
Near a holding voltage $V^0$, a gate pulse gives
$\delta H=\delta V(t)\widehat W$ with
$\widehat W=\partial H/\partial V$.
In a harmonic confinement model, a gate force $F_v\delta V$
has $\widehat W=-F_v x$, so
\[
 \Delta=\hbar\omega_x,\qquad
 \lambda(t)=\delta V(t)\bra\ell\widehat W\ket1
       =-F_v\delta V(t)\sqrt{\frac{\hbar}{2m^*\omega_x}}.
\]
Here $\Delta$ is the orbital gap and $F_v$ is force per voltage.
A uniform scalar shift would have no off-diagonal matrix element;
excitation occurs because this gate changes the shape or position
of the well. Analogous matrix elements couple to valley or charge
excitations when the corresponding potential changes allow them.

Set the energy of $\ket1$ to zero. For small pulses, neglect the
change of the gap and retain the weak coupling:
\begin{equation}
 H(t)=\Delta\ket\ell\bra\ell+
      \lambda(t)\ket\ell\bra1+\lambda^*(t)\ket1\bra\ell.
 \label{eq:calibration-leakage-model}
\end{equation}
The pulse begins and ends at the holding voltage, where these
states define the retained and excluded spaces. Starting in
$\ket1$, first-order time-dependent perturbation theory gives
\begin{equation}
 a_\ell(\tau)\simeq
 -\frac{i}{\hbar}\int_0^\tau
       \lambda(t)e^{i\Delta t/\hbar}\,dt.
 \label{eq:calibration-leakage-amplitude}
\end{equation}
The pulse component at angular frequency $\Delta/\hbar$ drives
the unwanted transition. For a constant weak coupling switched
on for time $\tau$,
\[
 |a_\ell(\tau)|^2\simeq
 \frac{4|\lambda|^2}{\Delta^2}
 \sin^2\!\left(\frac{\Delta\tau}{2\hbar}\right).
\]
\subsubsection{Comparing pulse shapes for the same intended action}
The leakage amplitude selects the pulse's Fourier component at
$\Delta/\hbar$. We now compare a square and a smooth pulse to see
which produces less excitation at that gap. Set $\omega_\ell=\Delta/\hbar$. A square pulse
$\lambda_{\mathrm{sq}}(t)=\lambda_0$ and a smooth pulse
\[
 \lambda_{\mathrm{sm}}(t)=
   2\lambda_0\sin^2(\pi t/\tau)
   =\lambda_0[1-\cos(2\pi t/\tau)],\qquad 0\le t\le\tau,
\]
have the same time integral. This comparison is appropriate when
the desired control and the unwanted coupling are both
proportional to the same waveform; otherwise the desired gate
must be matched using its own response.
Let $a=2\pi/\tau$. Direct integration of the three exponential
terms gives
\begin{align}
 \int_0^\tau\lambda_{\mathrm{sq}}e^{i\omega_\ell t}\,dt
   &=\lambda_0\frac{e^{i\omega_\ell\tau}-1}{i\omega_\ell},
       \notag\\
 \int_0^\tau\lambda_{\mathrm{sm}}e^{i\omega_\ell t}\,dt
   &=-\frac{a^2}{\omega_\ell^2-a^2}\,
       \lambda_0\frac{e^{i\omega_\ell\tau}-1}{i\omega_\ell}.
 \label{eq:calibration-smooth-pulse-spectrum}
\end{align}
Away from the removable singularities at $\omega_\ell=\pm a$,
the smooth-pulse probability is the square-pulse result times
$a^4/(\omega_\ell^2-a^2)^2$. For $\omega_\ell\tau\gg2\pi$,
this adds a suppression factor approximately
$(2\pi/\omega_\ell\tau)^4$. At $\omega_\ell=a$, the integral
must instead be evaluated by its finite limit; there is no
suppression guarantee near that frequency.

The smooth pulse and its first derivative vanish at both ends.
Those endpoint conditions remove the leading boundary terms
when integrating by parts and explain its smaller high-frequency
component. The comparison assumes perturbatively small leakage
and a fixed unwanted gap. If a voltage excursion brings an
excluded state close in energy, the gap must be included as
time dependent rather than using this fixed-gap formula.

\subsubsection{When the gap changes during a large ramp}
The preceding calculation assumed fixed unperturbed states and a
weak pulse. A large voltage excursion instead changes the
instantaneous eigenstates and their gaps. If the intended operation
follows one of those states,
the corresponding local adiabatic condition is
\begin{equation}
 \frac{\hbar|\bra{m(t)}\dot H(t)\ket{n(t)}|}
      {|E_m(t)-E_n(t)|^2}\ll1,\qquad m\ne n.
 \label{eq:calibration-adiabatic-condition}
\end{equation}
Apply it to the unwanted states being excluded, not automatically
to every pair of states in the computational space: the desired
gate may intentionally drive transitions within that space.
Use the measured excitation gaps to select candidate ramp times,
then recalibrate the gate angle for the changed waveform.
Longer pulses also accumulate more relaxation and phase noise,
so smoothness alone is not a performance criterion.

Changing the envelope also changes the intended gate.
For an electrically driven spin, replace the rectangular amplitude
$a_0$ of Section~\ref{sec:measurement-rabi} by a shaped
envelope $a(t)$. The corresponding Rabi angular frequency now varies:
\[
 \Omega(t)=\frac{g_{\mathrm e}\mu_Bb_\perp\chi_x a(t)}{2\hbar},
 \qquad
 H_{\mathrm{rot,q}}(t)
 =\frac{\hbar\Omega(t)}2(\cos\phi X-\sin\phi Y)
\]
at zero detuning and fixed phase $\phi$.
All these Hamiltonians are proportional to the same Pauli operator,
so the pulse produces angle
\[
 \vartheta=\int_0^{t_p}\Omega(t)\,dt.
\]
For a rectangular envelope this reduces to $\vartheta=\Omega_0t_p$;
a shaped pulse requires the full integral.
Reuse the Rabi or exchange measurement with the complete modified
pulse to restore its target angle. If detuning also varies, the
Hamiltonians at different times need not commute. The complete
propagator is then
\begin{equation}
 U(t_p)=\mathcal T\exp\!\left[-\frac{i}{\hbar}
                    \int_0^{t_p}H(t)\,dt\right],
 \label{eq:calibration-shaped-time-ordering}
\end{equation}
where $\mathcal T$ places later-time operators to the left. This is
the time-ordered evolution introduced in
Eq.~\eqref{eq:gate-time-evolution}; the integral cannot be replaced
by a scalar pulse area when the rotation axis changes in time.
Test both the intended gate and unwanted excitation before accepting
the new shape.
